\documentclass[conference]{IEEEtran}
\IEEEoverridecommandlockouts

\usepackage[utf8]{inputenc}
\usepackage[T1]{fontenc}
\usepackage{amsmath,amssymb,amsfonts}
\usepackage{graphicx}
\usepackage{booktabs}
\usepackage{listings}
\usepackage{xcolor}
\usepackage{balance}

\lstset{
  basicstyle=\ttfamily\scriptsize,
  columns=flexible,
  keepspaces=true,
  breaklines=true,
  breakatwhitespace=false,
  frame=single,
  numbers=none,
  xleftmargin=0pt,
  xrightmargin=0pt,
  showstringspaces=false,
  aboveskip=6pt,
  belowskip=6pt
}

\begin{document}

\title{Ordem de Criação versus Ordem de Término de Processos Filhos: Estudo Experimental com 
\texttt{fork()} e \texttt{wait()}}

\author{
\IEEEauthorblockN{Fábio Vinnicius Silva Sousa, Jonatas Andrade do Nascimento}
\IEEEauthorblockA{
Curso de Ciência da Computação,
Instituto Federal Campus Maracanaú
}
}

\maketitle

\begin{abstract}
Um processo pai que cria múltiplos processos filhos por meio de chamadas
sequenciais a \texttt{fork()} pode dar a falsa impressão de que esses filhos
também serão executados e terminarão na mesma ordem em que foram criados.
Este trabalho investiga experimentalmente essa hipótese. Foi implementado um
programa em C que cria $N$ processos filhos, cada um executando uma carga de
trabalho de duração aleatória e independente, e mede a ordem em que o
processo pai os recebe de volta via \texttt{wait()}. Os experimentos foram
repetidos centenas de vezes para diferentes valores de $N$ e comparados com o
modelo teórico de uma permutação aleatória uniforme. Os resultados mostram
que a hipótese deve ser \textbf{rejeitada}: a probabilidade de a ordem de
término coincidir com a ordem de criação cai rapidamente conforme $N$
aumenta, de forma consistente com $1/N!$, evidenciando que o escalonador do
sistema operacional não garante nenhuma relação entre ordem de criação e
ordem de término dos processos.
\end{abstract}

\begin{IEEEkeywords}
sistemas operacionais, fork, wait, waitpid, escalonamento de processos,
chamadas de sistema
\end{IEEEkeywords}

\section{Introdução}

Em sistemas operacionais do tipo Unix, um processo pai pode criar múltiplos
processos filhos por meio de chamadas sucessivas a \texttt{fork()}. Como os
filhos são criados sequencialmente, um raciocínio ingênuo poderia supor que
essa ordem se refletiria também na ordem em que os filhos terminam sua
execução e, consequentemente, na ordem em que o processo pai os recebe de
volta via \texttt{wait()}.

\textbf{Problema.} A ordem de criação dos processos filhos garante a ordem
em que eles serão executados e terminados?

\textbf{Hipótese.} A ordem em que os processos filhos são criados influencia
a ordem em que eles serão executados e terminados; consequentemente, ao
utilizar \texttt{wait()}, o processo pai tenderá a receber os filhos na
mesma ordem em que foram criados.

Esta hipótese é razoável à primeira vista, pois os filhos são de fato criados
em sequência e herdam código idêntico do pai. Entretanto, uma vez criado, cada
processo filho passa a ser uma entidade independente, sujeita ao escalonador
do sistema operacional e à duração do trabalho que executa, fatores que não
têm relação necessária com a ordem de criação. Este trabalho projeta e executa
um experimento controlado para testar essa hipótese.

\section{Metodologia}

\subsection{Desenho do experimento}

Foi implementado um programa em C (\texttt{experiment.c}) que, para um
parâmetro $N$ informado na linha de comando:

\begin{enumerate}
  \item cria $N$ processos filhos em laço sequencial usando \texttt{fork()},
        armazenando o PID de cada filho na ordem em que foi criado;
  \item cada filho executa uma carga de trabalho de duração aleatória e
        independente, sorteada uniformemente no intervalo $[10\,ms, 210\,ms]$
        via \texttt{usleep()}, e então encerra repassando seu índice de
        criação como código de saída (\texttt{\_exit(i)});
  \item o processo pai aguarda os $N$ filhos com \texttt{wait()} genérico
        (sem especificar qual PID esperar), registrando a ordem real em que
        os filhos terminam;
  \item ao final, o programa compara a sequência de PIDs de término com a
        sequência de PIDs de criação e reporta se a ordem foi preservada.
\end{enumerate}

\subsection{Cuidado metodológico: independência das amostras}

Uma armadilha identificada durante a implementação merece registro. Em uma
primeira versão, o gerador de números pseudoaleatórios era semeado
(\texttt{srand()}) uma única vez pelo processo pai, antes do laço de
\texttt{fork()}. Como o \texttt{fork()} duplica a memória do processo
(\emph{copy-on-write}), todos os filhos herdavam o \emph{mesmo} estado do
gerador e, portanto, sorteavam exatamente o mesmo tempo de trabalho, o que
invalidava a premissa de amostras independentes e enviesava artificialmente
o resultado a favor da hipótese (a ordem de término passava a depender
apenas da pequena vantagem de tempo de quem foi criado primeiro). A correção
foi semear o gerador \emph{dentro} de cada filho, após o \texttt{fork()},
combinando \texttt{clock\_gettime(CLOCK\_REALTIME)} com o PID do próprio
filho (Listagem~\ref{lst:seed}). Esse ponto ilustra como um detalhe de
implementação pode comprometer silenciosamente a validade de um experimento.

\begin{lstlisting}[
  caption={Semeadura independente por filho.},
  label={lst:seed}
]
if (pid == 0) {
    struct timespec ts;
    clock_gettime(CLOCK_REALTIME, &ts);
    srand((unsigned int)ts.tv_nsec ^ getpid());
    int tempo_us = (rand() % 200000) + 10000;
    usleep(tempo_us);
    _exit(i);
}
\end{lstlisting}

\subsection{Procedimento estatístico}

Para cada valor de $N \in \{2,3,4,5,10\}$, o programa foi executado
repetidamente (script \texttt{run\_experiments.sh}), e foi calculada a
fração de execuções em que a ordem de término coincidiu exatamente com a
ordem de criação. Esse valor observado foi comparado ao valor esperado sob
a hipótese nula de que a ordem de término é uma permutação uniformemente
aleatória das $N!$ permutações possíveis, isto é, probabilidade $1/N!$ de
coincidência exata por acaso.

\section{Resultados}

A Tabela~\ref{tab:resultados} apresenta, para cada $N$, o número de rodadas
executadas, a fração de rodadas em que a ordem de criação foi preservada e o
valor teórico $1/N!$ esperado sob aleatoriedade uniforme.

\begin{table}[htbp]
\caption{Fração de execuções em que a ordem de término coincidiu com a ordem de criação}
\label{tab:resultados}
\centering
\begin{tabular}{@{}ccccc@{}}
\toprule
$N$ & Rodadas & Ordem preservada & Observado & Esperado ($1/N!$) \\
\midrule
2  & 150 & 85  & 56{,}7\% & 50{,}0\% \\
3  & 150 & 18  & 12{,}0\% & 16{,}7\% \\
4  & 150 & 10  & 6{,}7\%  & 4{,}2\%  \\
5  & 150 & 0   & 0{,}0\%  & 0{,}8\%  \\
10 & 100 & 0   & 0{,}0\%  & $\approx$0{,}00003\% \\
\bottomrule
\end{tabular}
\end{table}

Os valores observados acompanham de perto a ordem de grandeza prevista pelo
modelo de permutação aleatória, especialmente para $N$ pequeno. Para
$N \geq 5$, a probabilidade teórica já é tão baixa que nenhuma das rodadas
executadas preservou a ordem original, como esperado. O padrão dominante é a
queda acentuada, praticamente fatorial, da taxa de preservação de ordem à
medida que o número de filhos cresce.

\emph{Nota:} os valores acima foram obtidos no ambiente de desenvolvimento
usado para este rascunho. Eles devem ser \textbf{regenerados na máquina da
equipe} antes da entrega, pois dependem do número de núcleos, da carga do
sistema e do escalonador do SO utilizado.

\section{Conclusões}

Os resultados experimentais rejeitam a hipótese original: a ordem de criação
dos processos filhos \emph{não} determina a ordem em que eles terminam nem a
ordem em que \texttt{wait()} os retorna ao processo pai. Um \texttt{wait()}
genérico simplesmente devolve o próximo filho que terminar, independente de
quando foi criado, e a probabilidade de a ordem coincidir com a criação
converge rapidamente para a de uma permutação aleatória conforme $N$ cresce.
Na prática, isso significa que qualquer código que dependa implicitamente da
ordem de criação dos filhos (por exemplo, para associar resultados a tarefas
específicas) é incorreto e deve, em vez disso, usar mecanismos explícitos de
identificação, como o valor de retorno de \texttt{fork()} armazenado pelo pai
combinado a \texttt{waitpid()} direcionado a um PID específico, quando a
ordem de atendimento precisar ser controlada.

\balance
\begin{thebibliography}{00}
\bibitem{silberschatz} A. Silberschatz, P. B. Galvin, and G. Gagne,
\emph{Operating System Concepts}, 10th ed. Hoboken, NJ, USA: Wiley, 2018.
\bibitem{stevens} W. R. Stevens and S. A. Rago,
\emph{Advanced Programming in the UNIX Environment}, 3rd ed. Boston, MA, USA: Addison-Wesley, 2013.
\bibitem{fork2} Linux manual page, \emph{fork(2)}. [Online]. Available: https://man7.org/linux/man-pages/man2/fork.2.html
\bibitem{wait2} Linux manual page, \emph{wait(2)}. [Online]. Available: https://man7.org/linux/man-pages/man2/wait.2.html
\end{thebibliography}

\end{document}
